% ================================================================
% RESEARCH CV
% Compile with XeLaTeX (Overleaf: Menu -> Compiler -> XeLaTeX).
% Requires Linux Biolinum O and Linux Libertine O (bundled with TeX Live / Overleaf).
% ================================================================
\documentclass[10pt,a4paper]{article}

% ---- Page, fonts and links --------------------------------------
\usepackage[left=0.44in,right=0.44in,top=0.68in,bottom=0.39in]{geometry}
\usepackage{fontspec} % XeLaTeX only
\defaultfontfeatures{Ligatures=TeX}
\setmainfont{lmroman10-regular.otf}[
  BoldFont=lmroman10-bold.otf,
  ItalicFont=lmroman10-italic.otf,
  BoldItalicFont=lmroman10-bolditalic.otf]
\setsansfont{Linux Biolinum O}
\newfontfamily\cvcontactfont{Linux Libertine O}
\usepackage{hyperref}
\hypersetup{colorlinks=true,allcolors=black,pdfborder={0 0 0},
  pdftitle={Yujian Cao -- Curriculum Vitae},pdfauthor={Yujian Cao}}
\usepackage{enumitem}
\usepackage{microtype}
\pagestyle{empty}
\raggedbottom\raggedright
\setlength{\parindent}{0pt}
\setlength{\parskip}{0pt}
\setlength{\emergencystretch}{1em}

% ---- Style controls: adjust these once for the entire CV --------
\newlength{\cvtitleindent}  \setlength{\cvtitleindent}{2pt}   % entry titles
\newlength{\cvdetailindent} \setlength{\cvdetailindent}{17pt} % detail lines
\newlength{\cvpubindent}    \setlength{\cvpubindent}{17pt}    % publication hanging indent
\newlength{\cvsectiongap}   \setlength{\cvsectiongap}{15pt}   % space above section headings
% Section: uppercase, bold Linux Biolinum, thin line underneath.
\newcommand{\cvsection}[1]{%
  \par\vspace{\cvsectiongap}%
  {\sffamily\bfseries\fontsize{12pt}{13pt}\selectfont\MakeUppercase{#1}}\par\nobreak
  \vspace{2.5pt}\hrule height 0.45pt\vspace{3.5pt}%
}

% Tighter rule spacing for the first heading on page 2.
\newcommand{\cvsectionpagetwo}[1]{%
  \par\vspace{\cvsectiongap}%
  {\sffamily\bfseries\fontsize{12pt}{13pt}\selectfont\MakeUppercase{#1}}\par\nobreak
  \vspace{-2.2pt}\hrule height 0.45pt\vspace{2.8pt}%
}

% A bold label at left; date aligned to the far right.
% LTL with italic subscript f.
\newcommand{\ltlf}{LTL\ensuremath{_{f}}}

\newcommand{\cvdated}[2]{%
  \par\noindent\hspace*{\cvtitleindent}\textbf{#1}\hfill\mbox{#2}\par\nobreak
}

% Education entry: bold institution, regular-weight location, date at right.
\newcommand{\cvedu}[3]{%
  \par\noindent\hspace*{\cvtitleindent}\textbf{#1}, #2\hfill\mbox{#3}\par\nobreak
}

% Indented details underneath a dated entry.
\newcommand{\cvdetail}[1]{%
  \noindent\hspace*{\cvdetailindent}#1\par\nobreak
}
\newcommand{\cventrygap}{\vspace{9pt}}

% Publications: fixed-width [1] numbering and hanging title indent.
\newcounter{cvpub}
\newcommand{\cvpublication}[3]{%
  \stepcounter{cvpub}%
  \par\noindent\hangindent=\cvpubindent\hangafter=1\hspace*{\cvtitleindent}%
  \makebox[\dimexpr\cvpubindent-\cvtitleindent\relax][l]{[\arabic{cvpub}]}#1\par\nobreak
  \noindent\hspace*{\cvpubindent}#2\par\nobreak
  \noindent\hspace*{\cvpubindent}#3\par\vspace{6pt}%
}
% Underlined placeholder; replace example.org with the actual PDF URL.
\newcommand{\cvpdf}[1]{\href{#1}{\underline{[PDF]}}}

% Awards and project bullets: short dash matching the reference.
\newenvironment{cvawards}{%
  \begin{itemize}[label={-},leftmargin=\cvdetailindent,labelsep=5pt,
    topsep=0pt,itemsep=3.5pt,parsep=0pt,partopsep=0pt]
}{\end{itemize}}
\newenvironment{cvbullets}{%
  \begin{itemize}[label={-},leftmargin=16pt,labelsep=5pt,
    topsep=1pt,itemsep=0pt,parsep=0pt,partopsep=0pt]
}{\end{itemize}}

% Project title and date; contents are written in cvbullets below.
\newcommand{\cvproject}[2]{\cvdated{#1}{#2}}
\newcommand{\cvprojectgap}{\vspace{12pt}}

% Industry/research appointment.
\newcommand{\cvexperience}[4]{%
  \cvedu{#1}{#2}{#3}%
  \cvdetail{\emph{#4}}%
}

\begin{document}
\fontsize{10pt}{14.15pt}\selectfont

{\sffamily\bfseries\fontsize{22pt}{25pt}\selectfont Yujian Cao}\par
\vspace{1pt}
{\cvcontactfont\fontsize{11pt}{13pt}\selectfont
  \href{mailto:yujian.cao@liverpool.ac.uk}{yujian.cao@liverpool.ac.uk}\enspace\textbar\enspace
  \href{https://yujian-cao.github.io/}{yujian-cao.github.io}\enspace\textbar\enspace
  \href{https://dblp.org/pid/160/1321.html}{\underline{[DBLP]}}
}\par\vspace{4pt}

\cvsection{Education}
\cvedu{University of Liverpool}{United Kingdom}{09/2025 -- Present}
\cvdetail{PhD in Computer Science}
\cvdetail{\emph{Supervisors: Dr. Shufang Zhu and Prof. Sven Schewe}}
\cventrygap

\cvedu{Hong Kong Baptist University}{Hong Kong SAR}{09/2022 -- 07/2024}
\cvdetail{MA in Artificial Intelligence and Digital Media}
\cventrygap

\cvedu{Anhui Agricultural University}{China}{09/2018 -- 06/2022}
\cvdetail{BEng in Computer Science and Technology}
\cventrygap

\cvedu{Anhui Normal University}{China}{09/2019 -- 12/2022}
\cvdetail{BA in Chinese Language and Literature}

\cvsection{Research Interests}
\noindent\hspace*{\cvtitleindent}Formal Methods, Reactive Synthesis, Linear Temporal Logic on Finite Traces (\ltlf{})\par

\cvsection{Experience}
\cvexperience{Festo}{China}{03/2025 -- 07/2025}
  {IT Technical Support Intern}
\cventrygap
\cvexperience{Sophia University}{Japan}{12/2024 -- 03/2025}
  {Visiting Student}
\cventrygap
\cvexperience{Beijing Normal--Hong Kong Baptist University}{China}{01/2024 -- 05/2024}
  {Research Assistant}

\cvsection{Scholarships \& Awards}
\begin{cvawards}
  \item 2026, University of Liverpool -- China Scholarship Council (CSC) Joint Scholarship
  \item 2024, Hong Kong Baptist University Merit Scholarship
  \item 2021, Anhui Agricultural University Scholarship
\end{cvawards}

\cvsection{Academic Service}
\noindent\hspace*{\cvtitleindent}Served as reviewer for ICASSP, ACL Rolling Review\par

\cvsection{Publications}
\cvpublication
  {Optimal \ltlf{} Synthesis.}
  {\textbf{Yujian Cao}, Sven Schewe, Qiyi Tang, Shufang Zhu}
  {In \emph{International Joint Conference on Artificial Intelligence} (IJCAI), 2026. \cvpdf{https://ijcai.org/proceedings/2026/0004.pdf}}

% \cvsection{Selected Research Projects}


% \cvproject{Efficient Foundation-Model Inference for Streaming Applications}{2023 -- Present}
% \begin{cvbullets}
%   \item Designed an adaptive routing algorithm that selects models according to query complexity and available resources, reducing inference cost while maintaining target accuracy on streaming benchmarks.
%   \item Developed an online optimization framework that combines predictive models with lightweight decision rules, supporting resource-aware execution across text and image-processing tasks, with deployment-time constraints for batch and interactive workloads in a shared cluster.
% \end{cvbullets}
% \cvprojectgap

% \cvproject{Unified Intermediate Representation for Data Query Languages}{2021 -- 2023}
% \begin{cvbullets}
%   \item Proposed a common intermediate representation that maps natural-language queries to structured programs in different execution environments.
%   \item Implemented a source-to-source translator supporting multiple graph and relational query formats, together with a reusable parser and evaluation toolkit.
%   \item Evaluated compositional generalization under domain transfer and low-resource settings, with consistent improvements over task-specific baseline systems, including unseen query structures and nested operations that are absent from the training set.
% \end{cvbullets}
% \cvprojectgap

% \cvproject{Automatic Performance Optimization for Software Systems}{2020 -- 2022}
% \begin{cvbullets}
%   \item Constructed a benchmark of performance regressions from public code repositories and trained a detection model using commit-level program features.
%   \item Studied compiler optimization selection and heterogeneous device scheduling, comparing learned policies against conventional heuristic approaches.
% \end{cvbullets}
% \cvprojectgap

% \cvproject{Representation Learning for Code Generation and Summarization}{2018 -- 2020}
% \begin{cvbullets}
%   \item Developed a contextual code representation model that uses historical changes to support automatic commit-message generation and semantic search.
%   \item Combined abstract syntax tree features with retrieval-based ranking to improve the relevance of generated explanations for source-code changes.
%   \item Explored graph-based information flow and syntax-aware attention for source-code summarization across several programming languages.
% \end{cvbullets}
% \cvprojectgap

% \cvproject{Task-Oriented Dialogue with External Knowledge Access}{2020}
% \begin{cvbullets}
%   \item Implemented a dialogue model that retrieves relevant passages from an external collection to answer questions beyond a structured task database.
%   \item Assessed response quality with automatic metrics and human evaluation, documenting error types and limitations under distribution shifts.
% \end{cvbullets}

\end{document}
